\section{The Sacrament of the Anointing of the Sick}
\label{sec:anointing}

Serious illness strikes the human person at a depth no purely clinical diagnosis can measure. It disrupts bodily strength, threatens cherished vocations, exposes creaturely fragility, awakens dread of death and judgment, and opens the heart to temptations of discouragement, bitterness, or despair. In the economy of the Incarnation, the Son of God did not abolish human suffering from a safe distance; He assumed mortal flesh, bore our infirmities (Matt~8:17; Is~53:4), wept before the tomb of Lazarus, agonized in Gethsemane, and turned bodily affliction on Calvary into the instrument of the world's redemption. Through the Sacrament of the Anointing of the Sick, the risen Lord draws near to the Christian weighed down by grave disease or the frailty of age, consecrating bodily weakness to His saving Cross, healing the lingering remnants of sin, and fortifying soul and body for either renewed earthly pilgrimage or peaceful passage into eternal glory.

\subsection{Scriptural Foundations: Mark 6 and the Apostolic Promulgation of James 5}
\label{subsec:anointing-scripture}

Throughout the Gospels, Jesus touches the sick with unmistakable compassion---laying His hands upon the fevered, touching the leper, healing the blind with clay and spittle---signifying that divine power enters human flesh through tangible means. When He dispatched the Twelve two by two on their Galilean mission, St.~Mark records a striking union of repentance and medicinal anointing:
\begin{quotation}
And going forth they preached that men should do penance: And they cast out many devils, and anointed with oil many that were sick, and healed them. (Mark~6:12--13; \latin{et ungebant oleo multos aegros, et sanabant})
\end{quotation}
As the Council of Trent (Session~XIV) and the \emph{Roman Catechism} (Part~II, \emph{On Extreme Unction}) teach, this apostolic practice did not originate in human ingenuity or natural pharmacy; Christ Himself commanded and intimated (\latin{insinuatum}) the sacred rite here, foreshadowing the sacrament that would receive its full New Covenant efficacy from His Passion and Resurrection.

The explicit apostolic promulgation of the sacrament stands in the Epistle of St.~James. Writing to the scattered Christian communities, the Apostle prescribes a determinate ecclesiastical rite for any believer struck by serious illness:
\begin{quotation}
Is any man sick among you? Let him bring in the priests of the church and let them pray over him, anointing him with oil in the name of the Lord. And the prayer of faith shall save the sick man. And the Lord shall raise him up: and if he be in sins, they shall be forgiven him. (Jas~5:14--15)
\end{quotation}
In the Latin Vulgate---\latin{Infirmatur quis in vobis? inducat presbyteros ecclesiae, et orent super eum, ungentes eum oleo in nomine Domini: et oratio fidei salvabit infirmum, et alleviabit eum Dominus: et si in peccatis sit, remittentur ei}---every essential element of a sacrament of the New Law is assembled in astonishing completeness:
\begin{enumerate}
  \item \textbf{The capable subject:} ``Is any man sick among you?'' (\latin{Infirmatur quis in vobis?})---a baptized member of the ecclesial community afflicted by serious bodily infirmity.
  \item \textbf{The proper minister:} ``Let him bring in the priests of the church'' (\latin{presbyteros ecclesiae})---not merely elderly members of the congregation or charismatic healers, but ordained presbyters of the Church.
  \item \textbf{The sensible matter and ritual act:} ``anointing him with oil in the name of the Lord'' (\latin{ungentes eum oleo}).
  \item \textbf{The deprecatory form:} ``let them pray over him\ldots\ and the prayer of faith'' (\latin{orent super eum\ldots\ oratio fidei}).
  \item \textbf{The supernatural grace and dual effect:} spiritual salvation and strengthening (``shall save the sick man. And the Lord shall raise him up''), extending even to the forgiveness of sins (``and if he be in sins, they shall be forgiven him''), alongside bodily alleviation as divine wisdom wills.
\end{enumerate}
As the \emph{Roman Catechism} observes directly from these words, when St.~James promises that ``if he be in sins, they shall be forgiven him,'' he attributes to this anointing with prayer the unmistakable nature and efficacy of a true sacrament instituted by Christ and promulgated to the Church through the Apostle.

\subsection{Patristic Testimony and Historical Development}
\label{subsec:anointing-patristics}

The Church's earliest liturgical and patristic monuments confirm that James~5:14--15 was received from the beginning as a sacred rite entrusted to the Church's priesthood rather than a passing charismatic gift of the apostolic generation. In fourth-century Egypt, the \emph{Euchologion} of St.~Serapion of Thmuis already preserves a formal episcopal blessing of oil for the sick, imploring God to send healing power from heaven for the cure of every bodily malady and every spiritual infirmity, and for the remission of sins. In the Greek East, St.~John Chrysostom (\emph{De sacerdotio}~III.6) cites James~5:14--15 expressly to show that the authority of Christian priests exceeds that of earthly parents: priests not only regenerate the soul in Baptism, but possess power to remit post-baptismal sins in the sick through anointing and prayer.

In the Latin West, the landmark canonical testimony is the decretal letter \emph{Si instituta} (Epistle~25, A.D.~416) written by Pope St.~Innocent~I to Decentius, Bishop of Gubbio. Answering Decentius's inquiry concerning the interpretation of James~5:14--15, Innocent~I states unequivocally that the Apostle's words refer to the faithful who are sick, that they are to be anointed with sacred oil blessed by the bishop, that not only presbyters but bishops themselves may celebrate the rite (since the bishop who has power to consecrate the oil certainly has power to visit and anoint the sick), and---decisively---that this anointing is a ``kind of sacrament'' (\latin{genus sacramenti}) which therefore cannot be administered to public penitents still excluded from sacramental communion. Innocent's reasoning proves that early fifth-century Rome already regarded the anointing of the sick as a true sacrament standing in the same ecclesial order as the Eucharist.

Three centuries later, St.~Bede the Venerable weaves Mark~6:13 and James~5:14--15 into a single apostolic tradition in his \emph{Exposition on the Gospel of Mark}:
\begin{quotation}
Where it is evident from the Apostles themselves, that it is an ancient custom of the holy Church that persons possessed or afflicted with any disease whatever, should be anointed with oil consecrated by priestly blessing. (\emph{In Marci evangelium expositio}~II.6, trans.\ Oxford \emph{Catena Aurea in Marcum}, 1842; Latin: \latin{Unde patet ab ipsis apostolis hunc sanctae Ecclesiae morem esse traditum, ut energumeni, vel alii quilibet aegroti ungantur oleo, pontificali benedictione consecrato})
\end{quotation}
While the 1842 Oxford rendering translates \latin{pontificali benedictione} as ``priestly blessing,'' Bede's Latin specifies episcopal consecration of the oil, echoing Innocent~I and the Roman liturgical tradition. In the same \emph{Catena Aurea} gloss on Mark~6:13, Pseudo-Chrysostom (Victor of Antioch) links the passage directly to St.~James and explains that oil signifies the mercy of God's unction, the healing of infirmity, and the enlightening of the heart worked through prayer, while Theophylact adds that oil denotes the grace of the Holy Spirit easing human labors and bestowing light and spiritual joy.

Over the centuries between the Carolingian reform and the high Middle Ages, Western pastoral practice underwent a significant narrowing of emphasis. Whereas the first millennium freely anointed the seriously sick in hope of both spiritual strengthening and bodily recovery, twelfth-century scholastic theologians began calling the rite \emph{Extreme Unction} (\latin{Extrema Unctio})---first simply to mark it as the final anointing after the anointings of Baptism, Confirmation, and Ordination, and gradually to treat it as the sacrament of the departing (\latin{sacramentum exeuntium}). The Council of Florence (\emph{Exsultate Deo}~\S14, 1439) reflected that medieval horizon when it stated that this fifth sacrament should be given only to a sick person whose death is feared (\latin{nisi infirmo, de cuius morte timetur}), with seven anointings upon the eyes, ears, nostrils, mouth, hands, feet, and loins.

Even at the height of the Tridentine era, however, the Church fought against the fatalistic habit of postponing the sacrament until agony or unconsciousness. In Session~XIV (1551), the Council of Trent solemnly defined against Protestant denials that Extreme Unction is a true and proper sacrament instituted by Christ and promulgated by St.~James, conferring the grace of the Holy Spirit, blotting out sins if any remain and the remnants of sin, comforting and strengthening the soul, and sometimes restoring bodily health when expedient for salvation. Fifteen years later, the \emph{Roman Catechism} (1566, Part~II, \emph{On Extreme Unction}) vigorously rebuked the grave sin of postponing the sacrament until all hope of recovery is gone and consciousness fades, teaching that receiving the anointing while mental faculties remain unimpaired enables the sick believer to elicit active faith and pious devotion and thus receive its graces more abundantly.

The Second Vatican Council brought that pastoral recovery to full maturity in \emph{Sacrosanctum Concilium}~73, teaching that ``Extreme Unction'' may also and more fittingly be called ``Anointing of the Sick,'' since it is not a sacrament exclusively for those at the point of death; rather, the fitting time to receive it has certainly arrived as soon as any of the faithful begins to be in danger of death from sickness or old age. On 30~November 1972, Pope St.~Paul~VI promulgated the apostolic constitution \emph{Sacram Unctionem Infirmorum}, renewing the Latin rite's matter, anointings, and essential form in fidelity to St.~James while preserving seamless doctrinal continuity with Florence and Trent.

\subsection{Scholastic Analysis, Matter, and Form}
\label{subsec:anointing-matter-form}

St.~Thomas Aquinas died before completing the \emph{Tertia Pars} of the \emph{Summa theologiae} beyond Penance, yet his mature plan in \emph{ST}~III, q.~65, aa.~1 and~3, together with his earlier commentary on Peter Lombard (\emph{In IV Sent.}, d.~23, subsequently compiled into \emph{Supplementum}, qq.~29--33), supplies the classic metaphysical anatomy of Anointing. Just as Baptism brings spiritual birth and Confirmation brings mature strength, so Penance restores spiritual life after mortal sin while Anointing of the Sick heals the lingering spiritual debility---the ``remnants of sin'' (\latin{reliquiae peccati})---that leaves the soul enfeebled in serious illness and unready for its passage to glory (\emph{ST}~III, q.~65, aa.~1, 3).

\paragraph{Matter (Remote and Proximate).}
The \textbf{remote matter} of the sacrament is blessed plant oil. Historically, Florence (\emph{Exsultate Deo}~\S14) and the \emph{Roman Catechism} specified olive oil consecrated by a bishop, praising oil's rich natural symbolism: it soothes pain, strengthens weary limbs, gives light, and refreshes fatigue, vividly mirroring the Spirit's healing grace within the afflicted soul. Exercising the Church's pastoral stewardship over changeable liturgical elements while preserving the substance instituted by Christ, Pope Paul~VI determined in \emph{Sacram Unctionem Infirmorum} (1972) that while olive oil remains preferred, another plant oil (\latin{oleum de plantis}) may validly be used according to local circumstances (CIC~1983, c.~999; CCEO~1990, c.~741). Ordinarily, the Oil of the Sick (\latin{Oleum Infirmorum}) is blessed by the diocesan bishop at the Chrism Mass on Holy Thursday; however, canon law also authorizes those equivalent in law to a diocesan bishop and, in case of necessity, any priest within the celebration of the sacrament itself to bless the oil (CIC~1983, c.~999). In the Eastern Catholic Churches, by ancient liturgical custom codified in CCEO~1990, c.~741, the priest administering the Mystery of Holy Unction (\latin{Euchelaion}) routinely blesses the oil during the celebration itself unless particular law reserves the blessing to the bishop.

The \textbf{proximate matter} is the ritual anointing of the sick person's body with the blessed oil. In the reformed Latin rite (\emph{Sacram Unctionem Infirmorum}; CIC~1983, c.~1000~\S1), the priest anoints the \textbf{forehead and the hands} of the sick person---the forehead signifying mind, faith, and inner perception, and the hands signifying human labor, touch, and active service---using his own hand (or, for a grave reason such as contagion, an instrument, c.~1000~\S2). In case of urgent necessity, a single anointing on the forehead, or even on another suitable part of the body, suffices for validity while reciting the entire form. In the Byzantine and other Eastern Catholic traditions, the Mystery of Holy Unction (\latin{Euchelaion}) unfolds through a solemn sequence of seven Epistle readings, seven Gospels, seven prayers, and seven anointings of the forehead, nostrils, cheeks, lips, breast, and hands---ideally celebrated by several priests (recalling the plural ``priests of the church'' in Jas~5:14), though validly administered by a single priest in ordinary pastoral need.

\paragraph{Form (Deprecatory Prayer).}
Unlike the indicative Latin forms of Baptism (``I baptize you\ldots'') and Penance (``I absolve you\ldots''), the sacramental form of Anointing has always been \textbf{deprecatory}---expressed as a solemn prayer of intercession. Anchored directly in St.~James's phrase ``let them pray over him\ldots\ and the prayer of faith shall save the sick man'' (Jas~5:14--15), the \emph{Roman Catechism} explains the deep fittingness of this prayerful mood: because the sacrament seeks not only the health of the soul but also, conditionally, the health of the body---which divine providence grants only when expedient for eternal salvation---the Church asks God in humble supplication rather than commanding bodily recovery.

In the reformed Latin rite promulgated by Pope Paul~VI (\emph{Sacram Unctionem Infirmorum}, 1972), the single essential form pronounced while anointing the forehead and hands is the Latin prayer \latin{Per istam sanctam unctionem et suam piissimam misericordiam adiuvet te Dominus gratia Spiritus Sancti, ut a peccatis liberatum te salvet atque propitius allevet}---praying that through this holy anointing and His most loving mercy the Lord may assist the sick believer with the grace of the Holy Spirit, freeing him from sin, saving him, and graciously raising him up. This concise formula gathers the exact language of James~5:14--15 into one unified prayer, replacing the older medieval sense-by-sense pardon petitions recorded in Florence (\S14) and the 1566 \emph{Roman Catechism}. In the Eastern Catholic Churches, the ancient deprecatory prayers likewise invoke the Father, Son, and Holy Spirit as the Heavenly Physician of souls and bodies, asking for the healing of every spiritual and bodily malady and the remission of sins.

\begin{threestable}{Ritual Dimension}{Classical Latin Tradition (Florence / 1566)}{Reformed Latin (1972 / CIC 1983) \& Eastern Catholic (\latin{Euchelaion} / CCEO 1990)}
Title and pastoral horizon & \emph{Extreme Unction} (\latin{Extrema Unctio}); restricted to danger of death (\emph{Exsultate Deo}~\S14), though \emph{Roman Catechism} warns against waiting for insensibility & \textbf{Latin:} \emph{Anointing of the Sick} (\emph{SC}~73); beginning of danger from grave illness or old age. \textbf{East:} \emph{Holy Unction} (\latin{Euchelaion}); grave sickness and, in some Eastern customs, solemn Lenten healing rites\\
Matter and blessing & Olive oil blessed by a bishop; anointing of the five organs of sense (plus feet, and formerly loins) & \textbf{Latin:} Olive oil or other plant oil, blessed by bishop or priest in necessity; forehead and hands (single anointing in necessity). \textbf{East:} Plant oil (often mixed with wine or water per ritual), routinely blessed by priest; multiple anointings\\
Essential form & Deprecatory prayer at each sense asking pardon for sins committed through that sense & \textbf{Latin:} Unified deprecatory prayer invoking the grace of the Holy Spirit to free from sin, save, and raise up (\emph{Sacram Unctionem}). \textbf{East:} Deprecatory prayers to the Heavenly Physician of souls and bodies\\
\end{threestable}

\subsection{Subject, Minister, and Repeatability}
\label{subsec:anointing-subject-minister}

\paragraph{The Capable Subject.}
The proper subject of the Anointing of the Sick is any living baptized Catholic who has attained the \textbf{use of reason} and begins to be in danger on account of \textbf{grave illness or old age} (CIC~1983, c.~1004~\S1; CCEO~1990, c.~738; \emph{Sacrosanctum Concilium}~73). Five theological and canonical boundaries define this eligibility:
\begin{enumerate}
  \item \textbf{Internal infirmity or frailty of age, not mere external danger:} The sacrament responds to sickness or bodily decline within the person---including serious chronic illness, frailty of advanced years, or imminent major surgery prompted by a dangerous illness---rather than external perils facing a healthy person, such as a soldier entering battle or a prisoner facing execution (\emph{Roman Catechism}, Part~II, \emph{On Extreme Unction}).
  \item \textbf{Use of reason:} Infants and persons who have permanently lacked the use of reason from birth are not capable subjects, because Anointing removes the remnants of personal post-baptismal sin and fortifies rational freedom against spiritual combat (\emph{Roman Catechism}, Part~II, \emph{On Extreme Unction}; \emph{Suppl.}~q.~32, a.~4). Children who have reached sufficient use of reason to be comforted by the sacrament may and should be anointed when gravely ill.
  \item \textbf{Living recipients only:} Sacraments are instituted exclusively for living wayfarers. A priest may never anoint a person known to be already dead; instead, he commends the departed soul to God's mercy in prayer. When prudent doubt remains whether death has actually occurred, whether the sick person has reached reason, or whether the illness is genuinely dangerous, the sacrament is conferred without scruple (CIC~1983, c.~1005).
  \item \textbf{Unconscious faithful and obstinate sinners:} The sacrament is to be conferred on sick persons who have lost consciousness or the use of their faculties if they would at least implicitly have asked for it while in possession of their senses (CIC~1983, c.~1006). Conversely, it is not to be conferred on those who obstinately persevere in manifest grave sin (CIC~1983, c.~1007), since deliberate impenitence places a direct obstacle (\latin{obex}) to grace.
  \item \textbf{Repeatability:} Because Anointing imprints no indelible character, it may be repeated whenever a sick person recovers after receiving it and falls gravely ill again, or whenever during the same prolonged illness the danger becomes more serious (CIC~1983, c.~1004~\S2; \emph{Roman Catechism}, Part~II, \emph{On Extreme Unction}).
\end{enumerate}

\paragraph{The Exclusive Priestly Minister.}
In unbroken fidelity to James~5:14 (``Let him bring in the priests of the church''), the Council of Trent (Session~XIV, ch.~3 and can.~4), and both modern Codes (CIC~1983, c.~1003~\S1; CCEO~1990, c.~739~\S1), \textbf{every priest and only a priest}---that is, a validly ordained bishop or presbyter---validly administers the Anointing of the Sick. Neither a deacon nor a layperson can ever validly confer this sacrament, as the Apostolic See has repeatedly reaffirmed. The theological reason is unmistakable: because Anointing extends instrumentally to the remission of sins and the remnants of sin, it belongs strictly to the ministerial priesthood endowed with the power of the keys in Christ the Head. Unlike Penance, however, Anointing does not require a special faculty of jurisdiction for validity: any validly ordained priest confers it validly, and in case of necessity any priest confers it licitly with at least the reasonably presumed permission of the pastor who has the care of souls (CIC~1983, c.~1003~\S2; CCEO~1990, c.~739~\S2).

\subsection{Metaphysical Effects, Order of Ends, and the Distinction from Viaticum}
\label{subsec:anointing-effects-viaticum}

Within the three-tier sacramental analysis, the Anointing of the Sick works a profound supernatural transformation of human frailty:
\begin{itemize}
  \item \textbf{\latin{Sacramentum tantum}:} The anointing of the gravely ill Christian with blessed plant oil accompanied by the Church's deprecatory prayer of faith.
  \item \textbf{\latin{Res et sacramentum}:} The \textbf{consecration and spiritual-bodily alignment} of the sick member to Christ suffering and victorious---a real, non-indelible configuration of the believer's infirmity to the redemptive Passion of the Head within His Mystical Body.
  \item \textbf{\latin{Res tantum}:} A rich fivefold medicinal grace of the Holy Spirit:
  \begin{enumerate}
    \item \emph{Strengthening, peace, and courage:} Quieting dread of death and judgment, dispelling spiritual gloom, and arming the believer against Satan's intensified temptations in the hour of weakness (\emph{Roman Catechism}, Part~II, \emph{On Extreme Unction}; CCC~1520).
    \item \emph{Union with Christ's Passion:} Elevating bodily suffering from passive endurance into an active spiritual sacrifice united to the Cross for the sanctification of the believer and the whole Church (Col~1:24; \emph{Lumen Gentium}~11; CCC~1521).
    \item \emph{Healing of the remnants of sin and venial sins:} Removing the spiritual languor, sluggishness, and weakened dispositions (\latin{reliquiae peccati}) left behind by past faults (\emph{ST}~III, q.~65, a.~1; \emph{Roman Catechism}, Part~II, \emph{On Extreme Unction}).
    \item \emph{Conditional remission of mortal sin:} Although Penance is the primary sacrament for mortal sin, if a gravely ill believer who has at least imperfect contrition is physically unable to confess (for instance, through sudden loss of speech or consciousness), Anointing infuses sanctifying grace and remits even mortal sins, fulfilling St.~James's promise: ``and if he be in sins, they shall be forgiven him'' (Jas~5:15; Trent, Sess.~XIV).
    \item \emph{Restoration of bodily health when expedient:} As Florence (\emph{Exsultate Deo}~\S14) and the \emph{Roman Catechism} affirm, the sacrament genuinely confers bodily alleviation and recovery insofar as it is expedient (\latin{in quantum autem expedit}) for the sick person's eternal salvation. When bodily cure is not granted, the grace is not frustrated; rather, it ripens into peace, fortitude, and readiness to meet the Lord.
  \end{enumerate}
\end{itemize}

Metaphysically, Anointing effects an \emph{accidental formal strengthening} of the soul through medicinal sacramental grace and a \emph{relational consecration} of the sick member to Christ crucified; it alters neither human nature nor the physical laws of mortality in a magical way, and imprints no indelible character.

Its \textbf{order of ends} exhibits the same harmonious hierarchy:
\begin{enumerate}
  \item \textbf{Primary proper end:} Spiritual healing, comfort, peace, and fortitude in grave illness or the frailty of old age, uniting the believer's suffering to the Passion of Christ.
  \item \textbf{Intrinsic secondary ends:} Eradication of the remnants of sin, remission of venial sins (and of mortal sins when a contrite sick person cannot confess), and restoration of bodily health if expedient for salvation.
  \item \textbf{Contingent and pastoral ends:} Testimony of Christian hope before family and medical caregivers, consolation of the Christian household, and spiritual readiness for recovery or surrender.
  \item \textbf{Ultimate end:} Peaceful consummation of the earthly pilgrimage and passage into the beatific glory of the risen Christ.
\end{enumerate}

\paragraph{Anointing of the Sick Distinguished from Viaticum.}
A crucial doctrinal and pastoral distinction guards the proper identity of Anointing: \textbf{Anointing of the Sick is not strictly the ``sacrament of the dying''---Viaticum is.} Popular speech often lumps Confession, Anointing, and final Communion together under the informal phrase ``Last Rites,'' blurring their distinct roles. Anointing of the Sick is the sacrament of \emph{grave illness and bodily frailty}, rightly celebrated early enough for the sick person to participate with alert faith, and repeatable if illness returns or worsens. \textbf{Viaticum} (\latin{viaticum}, ``provision for the journey'') is the \textbf{Holy Eucharist} received when a baptized Christian is actually in proximate danger of departing this life (CIC~1983, cc.~921--922; CCC~1524--1525).

Just as Baptism, Confirmation, and the Eucharist form the unified triad of Christian Initiation that launches the Christian journey, so \textbf{Penance, Anointing of the Sick, and Viaticum} form the continuous sacramental order that crowns it (\emph{Sacrosanctum Concilium}~74):

\begin{fourstable}{Sacrament in Crisis}{Proper Occasion}{Minister and Condition}{Role in the Final Passage}
Penance & Post-baptismal sin, whether at the onset of illness or in danger of death & Valid priest with faculties; \emph{any} priest in danger of death (CIC~c.~976) & Reconciles the soul with God and the Church, removing mortal guilt and restoring charity\\
Anointing of the Sick & Beginning of danger from grave illness or frailty of old age (repeatable) & Valid bishop or priest alone (CIC~c.~1003; CCEO~c.~739); living subject with reason & Heals sin's remnants, unites bodily affliction to Christ's Cross, strengthens soul and body as expedient\\
Viaticum (Holy Eucharist) & Proximate danger of departing this life from any cause (true sacrament of the dying) & Bishop, priest, deacon, or extraordinary minister (CIC~c.~911~\S2) & Feeds the dying Christian with the risen Christ as the pledge of eternal life and bodily resurrection (John~6:55 DR [6:54 Gk]; CCC~1524) for the passage to the Father\\
\end{fourstable}

In that ordered triptych of mercy, the Church reveals her tenderest motherhood. First she washes the wounded conscience in the blood of absolution; next she anoints the failing body with the oil of gladness and fortitude, sealing the sick member to the Cross of Christ; and finally, when the hour of departure truly arrives, she places upon the pilgrim's tongue the Living Bread come down from heaven, that he who eats Christ's Flesh and drinks His Blood may pass securely from death to everlasting life (John~6:55 DR [6:54 Gk]).
